% !PRINTED 70
\weberpage{70}

\begin{center}
{\large \S.\ 21.}

\vspace{0.5em}
{\large\gesperrt{Hauptsätze über Determinanten.}}
\end{center}

Aus dem Begriff der Determinante ergeben sich leicht die ersten Sätze, die für die Anwendung geeignet sind.

Wenn wir in dem Product [\S.\,20, (3)]
\begin{equation}
M'=\pm a_{\alpha_1}^{(1)}a_{\alpha_2}^{(2)}\cdots a_{\alpha_n}^{(n)}
\tag{1}
\end{equation}
die Factoren umstellen, so ändert sich sein Werth nicht. Wir können also die Factoren auch so anordnen, dass die unteren Indices in ihrer natürlichen Reihenfolge \(1,2,\ldots,n\) erscheinen. Dabei werden dann die oberen Indices in einer gewissen Weise permutirt erscheinen, also \(M'\) die Form erhalten:
\begin{equation}
\pm a_1^{(\beta_1)}a_2^{(\beta_2)}\cdots a_n^{(\beta_n)},
\tag{2}
\end{equation}
worin
\[
(\beta_1,\beta_2,\ldots,\beta_n)=\mathfrak B
\]
ebenso wie \((\alpha_1,\alpha_2,\ldots,\alpha_n)=\mathfrak A\) eine Anordnung der Ziffern \(1,2,\ldots,n\) bedeutet. Man kann die Anordnung \(\mathfrak B\) dadurch erhalten, dass man in den Factoren von \(M'\) die Transpositionen, die zu \(\mathfrak A\) geführt haben, von der letzten anfangend, rückgängig macht, um in der Reihe der unteren Indices wieder die ursprüngliche Anordnung zu erhalten. Die dabei sich ergebende Reihenfolge der oberen Indices ist dann die Anordnung \(\mathfrak B\). Es folgt daraus, dass \(\mathfrak B\) zur ersten oder zur zweiten Art gehört, je nachdem \(\mathfrak A\) zur ersten oder zur zweiten Art gehört, da beide durch die gleiche Anzahl von Transpositionen entstehen. Die Gesammtheit der \(\mathfrak B\) stellt ebenso wie die Gesammtheit der \(\mathfrak A\) alle Permutationen der \(n\) Elemente dar, da zwei verschiedene \(\mathfrak A\) niemals zu demselben \(\mathfrak B\) führen können. Damit ist bewiesen:

\gesperrt{I. Die Determinante \(\Delta\) kann auch dadurch gebildet werden, dass man in dem Hauptglied \(a_1^{(1)}a_2^{(2)}\cdots a_n^{(n)}\) die oberen Indices auf alle möglichen Arten permutirt, jedem der so gebildeten Producte das positive oder negative Zeichen giebt, je nachdem die angewandte Permutation zur ersten oder zweiten Art gehört, und dann die Summe aller dieser Producte nimmt.}
